% =============================================================================
% BEHAVIORAL IMPACT COCKPIT - PROJEKTPLAN
% BFE_019_ECHfP (Energie Schweiz für Private)
% Version 2.0
% =============================================================================

\documentclass[11pt,a4paper]{article}

% Packages
\usepackage[utf8]{inputenc}
\usepackage[T1]{fontenc}
\usepackage[ngerman]{babel}
\usepackage{geometry}
\usepackage{graphicx}
\usepackage{xcolor}
\usepackage{booktabs}
\usepackage{longtable}
\usepackage{tabularx}
\usepackage{multirow}
\usepackage{array}
\usepackage{fancyhdr}
\usepackage{titlesec}
\usepackage{enumitem}
\usepackage{tcolorbox}
\usepackage{tikz}
\usepackage{pgfgantt}
\usepackage{hyperref}

% Geometry
\geometry{margin=2.5cm, top=3cm, bottom=3cm}

% FehrAdvice Colors
\definecolor{fadarkblue}{RGB}{2,64,121}
\definecolor{falightblue}{RGB}{84,158,222}
\definecolor{fadarkgray}{RGB}{37,33,42}
\definecolor{falightgray}{RGB}{243,245,247}
\definecolor{faorange}{RGB}{222,157,62}
\definecolor{fared}{RGB}{192,80,77}
\definecolor{fagreen}{RGB}{126,189,172}

% Hyperref setup
\hypersetup{
    colorlinks=true,
    linkcolor=fadarkblue,
    urlcolor=fadarkblue,
    citecolor=fadarkblue
}

% Header/Footer
\pagestyle{fancy}
\fancyhf{}
\fancyhead[L]{\small\color{fadarkgray}Behavioral Impact Cockpit}
\fancyhead[R]{\small\color{fadarkgray}BFE-2026-001-COCKPIT}
\fancyfoot[C]{\small\color{fadarkgray}\thepage}
\fancyfoot[R]{\small\color{fadarkgray}Version 2.0}
\renewcommand{\headrulewidth}{0.5pt}
\renewcommand{\footrulewidth}{0.5pt}

% Section formatting
\titleformat{\section}
    {\Large\bfseries\color{fadarkblue}}
    {\thesection}{1em}{}
\titleformat{\subsection}
    {\large\bfseries\color{fadarkblue}}
    {\thesubsection}{1em}{}
\titleformat{\subsubsection}
    {\normalsize\bfseries\color{fadarkgray}}
    {\thesubsubsection}{1em}{}

% Custom boxes
\tcbset{
    fabox/.style={
        colback=falightgray,
        colframe=fadarkblue,
        boxrule=1pt,
        arc=3pt,
        left=10pt,
        right=10pt,
        top=10pt,
        bottom=10pt
    },
    warningbox/.style={
        colback=faorange!10,
        colframe=faorange,
        boxrule=1pt,
        arc=3pt,
        left=10pt,
        right=10pt,
        top=10pt,
        bottom=10pt
    },
    riskbox/.style={
        colback=fared!10,
        colframe=fared,
        boxrule=1pt,
        arc=3pt,
        left=10pt,
        right=10pt,
        top=10pt,
        bottom=10pt
    }
}

% New column types
\newcolumntype{L}[1]{>{\raggedright\arraybackslash}p{#1}}
\newcolumntype{C}[1]{>{\centering\arraybackslash}p{#1}}
\newcolumntype{R}[1]{>{\raggedleft\arraybackslash}p{#1}}

\begin{document}

% =============================================================================
% TITLE PAGE
% =============================================================================
\begin{titlepage}
    \centering
    \vspace*{2cm}

    {\color{fadarkblue}\rule{\textwidth}{2pt}}

    \vspace{1cm}

    {\Huge\bfseries\color{fadarkblue} Behavioral Impact Cockpit}

    \vspace{0.5cm}

    {\LARGE\color{fadarkgray} Vollständiger Projektplan}

    \vspace{1cm}

    {\color{fadarkblue}\rule{\textwidth}{2pt}}

    \vspace{2cm}

    \begin{tcolorbox}[fabox, width=0.8\textwidth]
        \centering
        \begin{tabular}{ll}
            \textbf{Projekt-ID:} & BFE-2026-001-COCKPIT \\
            \textbf{Parent-Projekt:} & BFE\_019\_ECHfP \\
            \textbf{Kunde:} & Bundesamt für Energie (BFE) \\
            \textbf{Timeline:} & 19.12.2025 -- 01.03.2026 \\
            \textbf{Version:} & 2.0 \\
            \textbf{Stand:} & 01.02.2026 \\
        \end{tabular}
    \end{tcolorbox}

    \vspace{2cm}

    {\large\color{fadarkgray}
    Dashboard zur Messung der Kampagnenwirkung\\
    über die KPIs Awareness (A), Willingness (W),\\
    Impact (I), Trust (T) + Gesamt-Score
    }

    \vfill

    {\color{fadarkgray}
    \textbf{FehrAdvice \& Partners AG}\\
    Klausstrasse 20, 8008 Zürich\\
    www.fehradvice.com
    }

\end{titlepage}

% =============================================================================
% TABLE OF CONTENTS
% =============================================================================
\tableofcontents
\newpage

% =============================================================================
% EXECUTIVE SUMMARY
% =============================================================================
\section{Executive Summary}

\begin{tcolorbox}[fabox]
Das \textbf{Behavioral Impact Cockpit} ist ein Dashboard zur Messung der Kampagnenwirkung von «Energie Schweiz für Private» (ECHfP). Es misst vier KPIs entlang der Behavioral Customer Journey:

\begin{itemize}[leftmargin=*]
    \item \textbf{A -- Awareness:} Bekanntheit der Kampagne (Gewicht: 25\%)
    \item \textbf{W -- Willingness:} Handlungsbereitschaft (Gewicht: 30\%)
    \item \textbf{I -- Impact:} Tatsächliche Verhaltensänderung (Gewicht: 30\%)
    \item \textbf{T -- Trust:} Vertrauen in BFE (Gewicht: 15\%)
\end{itemize}
\end{tcolorbox}

\subsection{Projektziele}

\begin{table}[h]
\centering
\rowcolors{2}{falightgray}{white}
\begin{tabular}{L{3cm}C{2.5cm}C{2.5cm}C{2.5cm}}
\toprule
\textbf{KPI} & \textbf{Baseline 2025} & \textbf{Ziel 2026} & \textbf{Stretch} \\
\midrule
Awareness & 45\% & 60\% & 70\% \\
Willingness & 25\% & 40\% & 50\% \\
Impact & 8\% & 15\% & 20\% \\
Trust & 55\% & 70\% & 80\% \\
\midrule
\textbf{Gesamt-Score} & -- & \textbf{≥ 0.45} & \textbf{0.55} \\
\bottomrule
\end{tabular}
\caption{KPI-Ziele für 2026}
\end{table}

\subsection{Projektkennzahlen}

\begin{table}[h]
\centering
\begin{tabular}{ll}
\toprule
\textbf{Kennzahl} & \textbf{Wert} \\
\midrule
Projektdauer & 11 Wochen \\
Budget & CHF 85'000 \\
Interne Stunden & 370h \\
Team-Grösse & 5 Personen (FehrAdvice) \\
Externe Partner & 3 (Intervista, Hakuna Matata, EVAI) \\
Arbeitspakete & 11 \\
Meilensteine & 6 \\
Review Gates & 4 \\
\bottomrule
\end{tabular}
\end{table}

\newpage

% =============================================================================
% MEILENSTEINE & TIMELINE
% =============================================================================
\section{Meilensteine \& Timeline}

\subsection{Meilenstein-Übersicht}

\begin{table}[h]
\centering
\rowcolors{2}{falightgray}{white}
\begin{tabular}{C{1cm}L{5cm}C{2.5cm}C{2cm}C{2cm}}
\toprule
\textbf{ID} & \textbf{Meilenstein} & \textbf{Datum} & \textbf{Status} & \textbf{Gate} \\
\midrule
MS1 & KPI-Architektur fertig & 19.12.2025 & \textcolor{fagreen}{✓ Erledigt} & -- \\
MS2 & Dashboard MVP fertig & 15.01.2026 & \textcolor{faorange}{◐ In Arbeit} & GATE-02 \\
MS3 & Fragebogen freigegeben & 23.01.2026 & ○ Ausstehend & GATE-01 \\
MS4 & Externe Daten integriert & 31.01.2026 & ○ Ausstehend & -- \\
MS5 & Textlösung fertig & 15.02.2026 & ○ Ausstehend & -- \\
MS6 & Cockpit übergeben & 01.03.2026 & ○ Ausstehend & GATE-04 \\
\bottomrule
\end{tabular}
\caption{Meilensteine mit Status und Review Gates}
\end{table}

\subsection{Gantt-Diagramm}

\begin{tcolorbox}[fabox]
\small
\begin{verbatim}
Dez 2025       |  Januar 2026                |  Februar 2026      |  März
15  19  22  26 |  02  08  15  22  29         |  05  12  19  26    |  01
---------------|----------------------------|--------------------|---------
AP1 ██         |                            |                    |
    └──────────|─AP2 ████████████           |                    |
               |              └──AP3 ████   |                    |
               |                    └───────|─AP4 ████████████████████████
               |                            |              └─AP9 ████
               |                            |                   └─AP8 ████
    └──────────|─AP5 ████████████           |                    |
               |              └──AP6 ███████|███                 |
               |              └──AP7 ███████|████████            |
AP11 ██████████|████████████████████████████|████████████████████|████
---------------|----------------------------|--------------------|---------
               |         ↑MS2      ↑MS3     |↑MS4      ↑MS5      |↑MS6
               |        15.01     23.01     |31.01    15.02      |01.03
\end{verbatim}
\end{tcolorbox}

\subsection{Kritischer Pfad}

\begin{tcolorbox}[warningbox]
\textbf{⚠ Kritischer Pfad (kein Puffer!)}

\vspace{0.3cm}
\textbf{AP1} → \textbf{AP2} → \textbf{AP3} → \textbf{AP4} → \textbf{AP9} → \textbf{AP8}

\vspace{0.3cm}
\textbf{Dauer:} 72 Tage | \textbf{Puffer:} 0 Tage

\vspace{0.3cm}
\textit{Jede Verzögerung auf diesem Pfad verzögert das Projektende direkt!}
\end{tcolorbox}

\newpage

% =============================================================================
% BUDGET & RESSOURCEN
% =============================================================================
\section{Budget \& Ressourcen}

\subsection{Budget-Übersicht}

\begin{table}[h]
\centering
\rowcolors{2}{falightgray}{white}
\begin{tabular}{L{5cm}R{2.5cm}R{2.5cm}R{2.5cm}}
\toprule
\textbf{Arbeitspaket} & \textbf{Intern (h)} & \textbf{Intern (CHF)} & \textbf{Extern (CHF)} \\
\midrule
AP1: KPI-Architektur & 40 & 8'000 & -- \\
AP2: Fragebogen & 60 & 12'000 & -- \\
AP3: Review \& Freigabe & 20 & 4'000 & -- \\
AP4: Briefing Intervista & 15 & 3'000 & 25'000 \\
AP5: Dashboard & 80 & 16'000 & -- \\
AP6: Integration ext. Daten & 30 & 6'000 & 2'000 \\
AP7: CMS/Textlösung & 25 & 5'000 & -- \\
AP8: Übergabe & 20 & 4'000 & -- \\
AP9: Auswertung & 40 & 8'000 & -- \\
AP10: Abstimmung HM & 10 & 2'000 & -- \\
AP11: Projektmanagement & 30 & 6'000 & -- \\
\midrule
\textbf{Subtotal} & \textbf{370} & \textbf{74'000} & \textbf{27'000} \\
Reserve (6\%) & -- & -- & 5'000 \\
\midrule
\textbf{Total} & \textbf{370} & \multicolumn{2}{r}{\textbf{CHF 85'000}} \\
\bottomrule
\end{tabular}
\caption{Budget pro Arbeitspaket}
\end{table}

\subsection{Ressourcen-Allokation}

\begin{table}[h]
\centering
\rowcolors{2}{falightgray}{white}
\begin{tabular}{L{3cm}L{5cm}C{2cm}C{2cm}}
\toprule
\textbf{Person} & \textbf{Rolle} & \textbf{Verfügb.} & \textbf{Stunden} \\
\midrule
Lucas & KPI-Architektur, Datenmodell & 40\% & 80 \\
Manuel & Dashboard, Auswertung & 50\% & 100 \\
Andrea & Fragebogen, Review, Partner & 30\% & 60 \\
Alexis & Projektleitung, Integration & 40\% & 80 \\
Isabella & Dashboard, Integration & 30\% & 60 \\
\midrule
\textbf{Total} & & & \textbf{380} \\
\bottomrule
\end{tabular}
\caption{Team-Ressourcen}
\end{table}

\newpage

% =============================================================================
% RISIKOMANAGEMENT
% =============================================================================
\section{Risikomanagement}

\subsection{Risiko-Übersicht}

\begin{table}[h]
\centering
\rowcolors{2}{falightgray}{white}
\begin{tabular}{C{0.8cm}L{4cm}C{1.8cm}C{1.8cm}C{1.5cm}L{2.5cm}}
\toprule
\textbf{ID} & \textbf{Risiko} & \textbf{W'keit} & \textbf{Auswirk.} & \textbf{Score} & \textbf{Verantw.} \\
\midrule
R1 & Verzögerung BFE-Freigabe & Hoch (60\%) & Hoch & \textcolor{fared}{\textbf{9}} & Andrea \\
R2 & Niedrige Rücklaufquote & Mittel (40\%) & Hoch & \textcolor{faorange}{\textbf{7}} & Andrea \\
R3 & Technische Probleme & Niedrig (20\%) & Mittel & \textcolor{fagreen}{\textbf{4}} & Manuel \\
R4 & Scope Creep durch BFE & Hoch (70\%) & Mittel & \textcolor{faorange}{\textbf{6}} & Alexis \\
R5 & Key Person Ausfall & Niedrig (15\%) & Hoch & \textcolor{fagreen}{\textbf{5}} & Alexis \\
\bottomrule
\end{tabular}
\caption{Risiko-Register}
\end{table}

\begin{tcolorbox}[fabox]
\textbf{Risiko-Matrix Legende:}
\begin{itemize}[leftmargin=*]
    \item \textcolor{fared}{\textbf{Score 8--10:}} Kritisch → Sofortmassnahmen erforderlich
    \item \textcolor{faorange}{\textbf{Score 5--7:}} Erhöht → Aktives Monitoring, wöchentlicher Review
    \item \textcolor{fagreen}{\textbf{Score 1--4:}} Akzeptabel → Review alle 2 Wochen
\end{itemize}
\end{tcolorbox}

\subsection{Risiko R1: Verzögerung BFE-Freigabe (Kritisch)}

\begin{tcolorbox}[riskbox]
\textbf{Beschreibung:} BFE benötigt mehr Zeit für Fragebogen-Review

\textbf{Auswirkung:} Verzögert Feldstart um 1--2 Wochen

\vspace{0.3cm}
\textbf{Präventive Massnahmen:}
\begin{itemize}[leftmargin=*]
    \item Vorab-Abstimmung mit BFE zu kritischen Fragen
    \item Fragebogen-Entwurf früh teilen (ab 10.01)
    \item Entscheidungstermine verbindlich vereinbaren
\end{itemize}

\textbf{Reaktive Massnahmen:}
\begin{itemize}[leftmargin=*]
    \item Parallele Arbeit an Dashboard fortsetzen
    \item Feldzeit verkürzen, Reminder verstärken
\end{itemize}

\textbf{Frühwarnung:} Keine Rückmeldung bis 20.01

\textbf{Eskalation bei:} Keine Freigabe bis 25.01
\end{tcolorbox}

\newpage

% =============================================================================
% ARBEITSPAKETE
% =============================================================================
\section{Arbeitspakete}

\subsection{AP1: KPI-Logik \& Architektur}

\begin{table}[h]
\centering
\begin{tabular}{L{4cm}L{8cm}}
\toprule
\textbf{Verantwortlich} & Lucas, Manuel \\
\textbf{Status} & \textcolor{fagreen}{✓ Abgeschlossen} \\
\textbf{Deadline} & 19.12.2025 \\
\textbf{Budget} & 40h / CHF 8'000 \\
\bottomrule
\end{tabular}
\end{table}

\textbf{Aufgaben:}
\begin{itemize}[leftmargin=*]
    \item[✓] AP1.1: Festlegen der Kennzahlen (A, W, I, T)
    \item[✓] AP1.2: Definieren der Berechnungslogik (Formeln, Gewichtungen)
    \item[✓] AP1.3: Erstellen des KPI-Datenmodells
    \item[✓] AP1.4: Dokumentation für Entwickler \& Datenteam
\end{itemize}

\textbf{Ergebnis:} Fertige KPI-Architektur und Formel-Set für das Cockpit

\subsection{AP2: Fragebogenerstellung}

\begin{table}[h]
\centering
\begin{tabular}{L{4cm}L{8cm}}
\toprule
\textbf{Verantwortlich} & Andrea, Lucas, Manuel \\
\textbf{Status} & \textcolor{faorange}{◐ In Arbeit} \\
\textbf{Deadline} & 15.01.2026 \\
\textbf{Budget} & 60h / CHF 12'000 \\
\bottomrule
\end{tabular}
\end{table}

\textbf{Aufgaben:}
\begin{itemize}[leftmargin=*]
    \item[◐] AP2.1: Entwicklung der Fragen für die Erstmessung
    \item[◐] AP2.2: Strukturierung nach Themen
    \item[◐] AP2.3: Definition Antwortskalen und Fragetypen
    \item[◐] AP2.4: KPI-Mapping sicherstellen
    \item[◐] AP2.5: Kommentierter Entwurf für Review
\end{itemize}

\textbf{Ergebnis:} Finaler Fragebogenentwurf zur internen Prüfung

\subsection{AP3: Review \& Freigabe}

\begin{table}[h]
\centering
\begin{tabular}{L{4cm}L{8cm}}
\toprule
\textbf{Verantwortlich} & Alexis, Andrea \\
\textbf{Status} & ○ Ausstehend \\
\textbf{Deadline} & 23.01.2026 \\
\textbf{Budget} & 20h / CHF 4'000 \\
\textbf{Risiko} & R1 (Kritisch) \\
\bottomrule
\end{tabular}
\end{table}

\textbf{Aufgaben:}
\begin{itemize}[leftmargin=*]
    \item[○] AP3.1: Gemeinsame Durchsicht (FAP, c-rk, BFE)
    \item[○] AP3.2: Prüfung Verständlichkeit und Relevanz
    \item[○] AP3.3: Einarbeitung Änderungswünsche
    \item[○] AP3.4: BFE-Freigabe einholen
\end{itemize}

\textbf{Ergebnis:} Offiziell freigegebener Fragebogen

\subsection{AP4: Briefing Intervista}

\begin{table}[h]
\centering
\begin{tabular}{L{4cm}L{8cm}}
\toprule
\textbf{Verantwortlich} & Andrea, Alexis \\
\textbf{Status} & ○ Ausstehend \\
\textbf{Deadline} & 15.03.2026 \\
\textbf{Budget} & 15h / CHF 28'000 (inkl. Intervista) \\
\textbf{Risiko} & R2 (Erhöht) \\
\bottomrule
\end{tabular}
\end{table}

\subsection{AP5: Dashboard-Entwicklung}

\begin{table}[h]
\centering
\begin{tabular}{L{4cm}L{8cm}}
\toprule
\textbf{Verantwortlich} & Alexis, Manuel, Isabella \\
\textbf{Status} & \textcolor{faorange}{◐ In Arbeit} \\
\textbf{Deadline} & 15.01.2026 \\
\textbf{Budget} & 80h / CHF 16'000 \\
\bottomrule
\end{tabular}
\end{table}

\textbf{Aufgaben:}
\begin{itemize}[leftmargin=*]
    \item[✓] AP5.1: MVP-Frontend erstellen
    \item[✓] AP5.2: KPI-Grafiken aufbauen
    \item[✓] AP5.3: Datenschnittstellen implementieren
    \item[✓] AP5.4: Filter gestalten
    \item[✓] AP5.5: Exportfunktionen sicherstellen
\end{itemize}

\textbf{Ergebnis:} Funktionsfähiger Dashboard-Prototyp (MVP)

\newpage

% =============================================================================
% GOVERNANCE & KOMMUNIKATION
% =============================================================================
\section{Governance \& Kommunikation}

\subsection{Rollen \& Verantwortlichkeiten}

\begin{table}[h]
\centering
\rowcolors{2}{falightgray}{white}
\begin{tabular}{L{3.5cm}L{4cm}L{5cm}}
\toprule
\textbf{Rolle} & \textbf{Person} & \textbf{Befugnisse} \\
\midrule
Projektleiter & Alexis & Entscheidungen bis CHF 5'000, Terminverschiebungen bis 5 Tage \\
Stellvertreter & Andrea & Wie Projektleiter bei Abwesenheit \\
Steering Committee & Andrea (FAP), BFE & Entscheidungen > CHF 5'000, Scope-Änderungen \\
\bottomrule
\end{tabular}
\caption{Governance-Struktur}
\end{table}

\subsection{Eskalationspfad}

\begin{table}[h]
\centering
\rowcolors{2}{falightgray}{white}
\begin{tabular}{C{1.5cm}L{4.5cm}L{4cm}C{2cm}}
\toprule
\textbf{Level} & \textbf{Trigger} & \textbf{Entscheider} & \textbf{Reaktion} \\
\midrule
1 & Verzögerung < 5 Tage, < CHF 2'000 & Projektleiter & 24h \\
2 & Verzögerung 5--10 Tage, CHF 2--5'000 & PL + Partner & 48h \\
3 & Verzögerung > 10 Tage, > CHF 5'000 & Steering Committee & 1 Woche \\
\bottomrule
\end{tabular}
\caption{Eskalationsstufen}
\end{table}

\subsection{Kommunikationsplan}

\begin{table}[h]
\centering
\rowcolors{2}{falightgray}{white}
\begin{tabular}{L{4cm}L{3cm}L{3.5cm}L{2.5cm}}
\toprule
\textbf{Meeting} & \textbf{Rhythmus} & \textbf{Teilnehmer} & \textbf{Output} \\
\midrule
Daily Standup & Mo, Mi, Fr 09:00 & Projektteam & Teams-Update \\
Weekly Review & Fr 14:00 & Team + Andrea & Status-Update \\
BFE Status & Alle 2 Wochen & Alexis, Andrea, BFE & Minutes \\
Meilenstein-Report & Bei MS & Alexis → BFE & PDF-Report \\
\bottomrule
\end{tabular}
\caption{Meeting-Übersicht}
\end{table}

\newpage

% =============================================================================
% QUALITÄTSSICHERUNG
% =============================================================================
\section{Qualitätssicherung}

\subsection{Review Gates}

\begin{table}[h]
\centering
\rowcolors{2}{falightgray}{white}
\begin{tabular}{C{1.5cm}L{4cm}C{2.5cm}L{5cm}}
\toprule
\textbf{Gate} & \textbf{Name} & \textbf{Datum} & \textbf{Kriterien} \\
\midrule
GATE-01 & Fragebogen-Freigabe & 23.01.2026 & BFE-Freigabe, Pretest < 5\% Probleme \\
GATE-02 & Dashboard Go-Live & 15.01.2026 & Alle Features funktional, Performance < 3s \\
GATE-03 & Datenfreigabe & 08.02.2026 & n ≥ 1'500, Qualitäts-Checks bestanden \\
GATE-04 & Projektabnahme & 01.03.2026 & Alle Deliverables, Dokumentation vollständig \\
\bottomrule
\end{tabular}
\caption{Review Gates}
\end{table}

\subsection{Qualitätskriterien Fragebogen}

\begin{itemize}[leftmargin=*]
    \item \textbf{Verständlichkeit:} Pretest mit 10 Personen, < 5\% Probleme
    \item \textbf{Bearbeitungszeit:} ≤ 15 Minuten (Median)
    \item \textbf{KPI-Abdeckung:} Jede Frage ist einem KPI zugeordnet
    \item \textbf{EBF-Integration:} ≥ 8 von 10 Dimensionen abgedeckt
    \item \textbf{Sprachversionen:} Native Speaker Review DE/FR/IT
\end{itemize}

\subsection{Qualitätskriterien Dashboard}

\begin{itemize}[leftmargin=*]
    \item \textbf{Funktionalität:} 100\% der Features funktional (11-Punkte-Checkliste)
    \item \textbf{Performance:} Ladezeit < 3 Sekunden
    \item \textbf{Usability:} 3 BFE-Nutzer können ohne Anleitung navigieren
    \item \textbf{Datenqualität:} 0\% Abweichung bei KPI-Berechnung
\end{itemize}

\subsection{Qualitätskriterien Daten}

\begin{itemize}[leftmargin=*]
    \item \textbf{Stichprobengrösse:} n ≥ 1'500 vollständige Interviews
    \item \textbf{Repräsentativität:} Quotiert nach Region, Alter, Gebäudetyp
    \item \textbf{Datenqualität:} < 5\% Ausschluss (Speeders, Straightliners)
    \item \textbf{Vollständigkeit:} < 3\% Missing Values bei Pflichtfragen
\end{itemize}

\newpage

% =============================================================================
% TEAM
% =============================================================================
\section{Team \& Partner}

\subsection{FehrAdvice Team}

\begin{table}[h]
\centering
\rowcolors{2}{falightgray}{white}
\begin{tabular}{L{3cm}L{6cm}C{2.5cm}}
\toprule
\textbf{Name} & \textbf{Rolle} & \textbf{Schwerpunkt-APs} \\
\midrule
Lucas & KPI-Architektur, Datenmodell, Auswertung & AP1, AP2, AP9 \\
Manuel & Dashboard-Entwicklung, Auswertung & AP1, AP5, AP9 \\
Andrea & Fragebogen, Review, Übergabe, Partner & AP2, AP3, AP8 \\
Alexis & Projektleitung, Dashboard, Integration & AP3, AP5--7, AP11 \\
Isabella & Dashboard-Entwicklung, Integration & AP5, AP6 \\
\bottomrule
\end{tabular}
\caption{FehrAdvice Projektteam}
\end{table}

\subsection{Externe Partner}

\begin{table}[h]
\centering
\rowcolors{2}{falightgray}{white}
\begin{tabular}{L{4cm}L{8cm}}
\toprule
\textbf{Partner} & \textbf{Rolle} \\
\midrule
Intervista & Feldarbeit, Online-Befragung (n=1'500+) \\
Hakuna Matata & Abstimmung Auswertungsplan \\
EVAI & Social-Listening-Tool (iFrame-Integration) \\
\bottomrule
\end{tabular}
\caption{Externe Partner}
\end{table}

\subsection{Kunde}

\begin{table}[h]
\centering
\rowcolors{2}{falightgray}{white}
\begin{tabular}{L{4cm}L{8cm}}
\toprule
\textbf{Stakeholder} & \textbf{Rolle} \\
\midrule
BFE & Auftraggeber, Freigabe, Abnahme \\
c-rk & Review \\
\bottomrule
\end{tabular}
\caption{Kundenseite}
\end{table}

\newpage

% =============================================================================
% EBF INTEGRATION
% =============================================================================
\section{EBF-Integration}

Das Cockpit basiert auf dem \textbf{Evidence-Based Framework (EBF)} und misst die Kampagnenwirkung entlang der \textbf{Behavioral Customer Journey (BCJ)}.

\subsection{KPI zu 10C Mapping}

\begin{table}[h]
\centering
\rowcolors{2}{falightgray}{white}
\begin{tabular}{C{1.5cm}L{3cm}C{2.5cm}L{5cm}}
\toprule
\textbf{KPI} & \textbf{10C-Dimension} & \textbf{EBF-Variable} & \textbf{Messung} \\
\midrule
A & AWARE & A(·) & Bekanntheit Kampagne, Förderungen \\
W & READY & W\textsubscript{base}, θ & Handlungsbereitschaft, Schwelle \\
I & STAGE & BCJ Phase φ & Verhaltensänderung \\
T & WHAT & u\textsubscript{S} & Vertrauen, Social Utility \\
\bottomrule
\end{tabular}
\caption{Mapping KPIs zu 10C-Dimensionen}
\end{table}

\subsection{BCJ-Integration}

\begin{table}[h]
\centering
\rowcolors{2}{falightgray}{white}
\begin{tabular}{L{3cm}C{2cm}L{6cm}}
\toprule
\textbf{BCJ-Phase} & \textbf{Primär-KPI} & \textbf{Interventions-Fokus} \\
\midrule
Unaware & A & Reichweite erhöhen \\
Aware & A, T & Vertrauen aufbauen \\
Considering & W & Barrieren abbauen \\
Intending & W, I & Commitment stärken \\
Acting & I & Umsetzung unterstützen \\
Maintaining & I, T & Zufriedenheit sichern \\
\bottomrule
\end{tabular}
\caption{BCJ-Phasen und KPI-Zuordnung}
\end{table}

\begin{tcolorbox}[fabox]
\textbf{Gesamt-Score Formel:}

\vspace{0.3cm}
\centering
\Large
$G = 0.25 \cdot A + 0.30 \cdot W + 0.30 \cdot I + 0.15 \cdot T$

\vspace{0.3cm}
\normalsize
\raggedright
Die Gewichtung reflektiert die strategische Priorität: \textbf{Willingness} und \textbf{Impact} haben den höchsten Einfluss (je 30\%), da tatsächliche Verhaltensänderung das Hauptziel der Kampagne ist.
\end{tcolorbox}

\newpage

% =============================================================================
% ANHANG: CHECKLISTEN
% =============================================================================
\section{Anhang: Checklisten}

\subsection{Fragebogen-Freigabe Checkliste}

\begin{itemize}[leftmargin=*]
    \item[$\square$] Pretest durchgeführt (n=10)
    \item[$\square$] Bearbeitungszeit ≤ 15 min
    \item[$\square$] Alle Fragen KPI-gemapped
    \item[$\square$] Übersetzungen geprüft (DE/FR/IT)
    \item[$\square$] BFE-Review abgeschlossen
    \item[$\square$] Freigabe-Mail von BFE erhalten
\end{itemize}

\subsection{Dashboard Go-Live Checkliste}

\begin{itemize}[leftmargin=*]
    \item[$\square$] KPI A angezeigt
    \item[$\square$] KPI W angezeigt
    \item[$\square$] KPI I angezeigt
    \item[$\square$] KPI T angezeigt
    \item[$\square$] Gesamt-Score berechnet
    \item[$\square$] Filter Zeit funktional
    \item[$\square$] Filter Zielgruppe funktional
    \item[$\square$] Filter Region funktional
    \item[$\square$] Export PDF funktional
    \item[$\square$] Export PowerPoint funktional
    \item[$\square$] Export CSV funktional
\end{itemize}

\subsection{Projektabschluss Checkliste}

\begin{itemize}[leftmargin=*]
    \item[$\square$] Alle Deliverables übergeben
    \item[$\square$] User Guide erstellt
    \item[$\square$] Übergabepräsentation durchgeführt
    \item[$\square$] Abnahmeprotokoll unterschrieben
    \item[$\square$] Lessons Learned dokumentiert
    \item[$\square$] Archivierung abgeschlossen
\end{itemize}

\vfill

\begin{center}
{\color{fadarkgray}\rule{0.5\textwidth}{0.5pt}}

\vspace{0.5cm}

{\small\color{fadarkgray}
\textbf{Dokument erstellt am:} 01.02.2026\\
\textbf{FehrAdvice \& Partners AG}\\
Behavioral Impact Cockpit -- Projektplan v2.0
}
\end{center}

\end{document}
